\section{Reducing the adverse word to two supports}
\label{sec:two-support-census}

The closed-box program can answer a more informative question than the value of its highest peak. Set the trial fraction below the known construction and ask for the identity of \emph{every} support that could clear it. If very few survive, the nonlinear marked-cell restrictions can be brought back on those supports individually. This reverses the order of work from adding a single global inequality and improves on the price bound even when its maximizing rate vector changes discontinuously.

\subsection{A change of relative phase}

The decisive word is \(W=2549776587809433532\). It arose from the sampling and recombination of labeled 32-leaf halves, followed by a change of coordinates on one half. There is a useful control with an exact geometric interpretation. Start at word \(7382425669205845932\). Keep all even physical leaves fixed. In the ranking of the odd leaves \(2y+1\), replace each \(y\) by \(11y+6\bmod 32\) and transport the rankings with the leaves. Keep the even half preferred at the root. Reading off preferences from the resulting rankings gives \(W\).

Multiplication by \(11\) permutes the odd directions on the 32-cycle and its dyadic quotients. Translation permutes the residue nodes. Consequently the two odd-half problems have identical gain and allowance data after relabeling, including their preferred children. Their relative positions with respect to the \emph{fixed} even half have changed. In the two-sector relaxation with cap \(12/5\) and at least two preferred root leaves in its missing-step sector, the original word gives \(32/187\) and the transported word gives \(58/567\). This large drop is a root effect. It shows why optimizing each half and accepting the first interleaving of their coordinates leaves an important decision uncontrolled.

\subsection{All supports above a threshold}

Fix \(z=10118/100000\), corresponding to \(10.118\%\). Two cap enclosures will be useful,
\[
  C_1=23921/10000,\qquad C_2=239213/100000.
\]
In each case a vertex for fixed structural support has 26 leaves of rate \(C_j\) and one leaf of rate \(64-26C_j\). Structural support here means all leaves regarded as occupied when calculating edge availability. A leaf in that support may carry zero mass at an optimizing vertex of the closure. The enumeration keeps these leaves as occupied. This detail is essential when an argument is to exclude all supports with positive rates as well as their independently optimized boundary vertices.

Let \(H_B(A,k,j)\) be the best contribution of the individual half mask \(A\) on side \(B\) when \(k\) full slots and \(j\) partial slots are placed there. Include the root \emph{charge} in the odd half and temporarily credit zero root gain. The best possible completion in the opposite half is
\[
       F_{\overline B}(26-k,1-j)=
       \max_{A'} H_{\overline B}(A',26-k,1-j).
\]
A necessary condition for a positive surplus on \(A\) is
\[
        H_B(A,k,j)+F_{\overline B}(26-k,1-j)>0.
\]
Apply this filter to all 3,263,443 legal masks on each side, separately for each slot split, then form the actual cross products of the two surviving lists and check their sums. This last check prevents two half masks found using different best completions from being credited with a positive joint surplus by mistake.

There are two positive full masks. Both use the same slot split at both cap enclosures. The physical even half carries 21 full slots and the partial slot. The odd half carries the five remaining full slots. Exactly two even masks and one odd mask clear the opposing-envelope filter at that split. Both compatible pairings have positive surplus. In the distinct sector where a root step is missing, the full-credit calculation examined 42,384 compatible support--direction pairs, comprising 39,058 distinct supports with odd mass and at least two preferred leaves. Even its best surplus was negative. The cap proof and the preferred-root mass constraint justify these relaxations within the stated marked-cell model.

\begin{center}\small
\begin{tabular}{@{}lrr@{}}
\toprule
Greatest \(P-zQ\) in the indicated class & \(C_1\) & \(C_2\)\\
\midrule
Missing root step, credited & \(-26.536285548\) & \(-26.5321517244\)\\
First surviving mask, root uncredited & \(0.378878364\) & \(0.3815630892\)\\
Second surviving mask, root uncredited & \(2.315139788\) & \(2.3178487964\)\\
Best other mask, root uncredited & \(-3.191590756\) & \(-3.1863520468\)\\
\bottomrule
\end{tabular}
\end{center}

All decimals in this table are exact. Integer arithmetic multiplies \(P-zQ\) by the denominator of the cap and by \(100000\). The first and second surviving masks are \(1230983613671479063\) and \(1518083709116552983\), with bit \(t\) denoting residue \(t\). Each has 28 leaves. The same distinct support \(1518085900086613781\) maximizes the uncredited envelope after removing the two winners at either cap. Thus the separation from all other supports is visible as a finite signed margin.

\subsection{A second computation after removing the winners}

An independent calculation avoids the opposing-envelope filter. It enumerates legal quarter masks by taking their preferred-child closure, joins the quarters into legal halves, and computes gains using the forbidden translates for the quarter pair. For each half and slot split retain the three highest scoring, distinct masks, with any consistent way of breaking ties. Join the resulting short lists and forbid the two specified full masks. This computes the exact best remaining surplus.

To see why three suffice, hold an opposite mask fixed. At most two masks on this side could complete it to one of the two forbidden full supports. Thus at least one of the three highest entries is compatible with the exclusion. If a surviving optimum uses an entry below the top three, replace that entry without decreasing the surplus. Having made the replacement on the first side, make the same argument on the second with the new first side held fixed. The resulting admissible pair lies in the cross product of the two length-three lists. If a half has fewer than three candidates they are all present. This argument includes ties and requires no assumption that the maximizers are unique. The fourth row of the table comes from this separate optimization. Both programs are replayed by \texttt{evidence/rootmark/near.py}.

What remains is a two-support calculation with arbitrary real rates. In both supports the preferred singleton of the sparse 32-leaf half is residue 27 and the other low-cost adjustment is residue 46. The counts of recoverable directions agree in the two saturated ledgers; the first carries an additional eight units of allowance before multiplication by the common productive rate. The next argument determines how much room any redistribution of mass has once the sparse singleton must satisfy the original marked-cell constraints.
